% !TeX program = xelatex
\documentclass[UTF8,a4paper,zihao=-4]{ctexart}

\usepackage{geometry}
\usepackage{fontspec}
\usepackage{graphicx}
\usepackage{fancyhdr}
\usepackage{titlesec}
\usepackage{setspace}
\usepackage{enumitem}
\usepackage{booktabs}
\usepackage{array}
\usepackage{longtable}
\usepackage{amsmath,amssymb,bm}
\usepackage{xcolor}
\usepackage{caption}
\usepackage{hyperref}

\geometry{
  top=2.54cm,
  bottom=2.54cm,
  left=3.17cm,
  right=3.17cm,
  headheight=1.60cm,
  headsep=0.60cm,
  footskip=1.75cm
}

\linespread{1.5}
\setlength{\parindent}{2em}
\setlength{\parskip}{0pt}
\hypersetup{hidelinks}

\titleformat{\section}[block]{\centering\bfseries\zihao{3}}{}{0pt}{}
\titlespacing*{\section}{0pt}{1.2em}{0.8em}
\setlist{nosep,leftmargin=2em}
\captionsetup{font=small,labelfont=bf}

\newcommand{\ImagePlaceholder}[2]{%
  \begin{figure}[htbp]
    \centering
    \fbox{\parbox[c][#2][c]{0.86\textwidth}{\centering #1}}
    \caption{#1}
  \end{figure}
}

\pagestyle{fancy}
\fancyhf{}
\fancyhead[C]{基于RDK X5的智能无影灯系统}
\fancyfoot[C]{\thepage}
\renewcommand{\headrulewidth}{0pt}
\renewcommand{\footrulewidth}{0pt}

\begin{document}

\begin{center}
{\fontsize{18pt}{22pt}\selectfont 基于RDK X5的智能无影灯系统}
\end{center}

\section*{摘要}

本作品面向阅读、书写、实验操作等近距离作业场景，设计并实现了一套基于RDK X5的智能无影灯系统。系统以RDK X5为嵌入式核心平台，在Ubuntu 22.04与ROS2 Humble环境下融合Intel RealSense D435视觉感知、板端离线语音交互、OpenMANIPULATOR-X机械臂定向调节以及基于SPI驱动的WS2812B灯光控制，形成“感知、决策、执行、反馈”一体化闭环。视觉侧将手部、阴影与姿态状态识别部署在RDK X5 BPU上，实现对桌面作业区域遮挡阴影的实时感知与补光目标估计；语音侧支持跟踪开关、灯光启停、亮度升降、最大/最小/中等亮度以及冷暖光切换等离线指令；协调侧通过ROS2话题机制完成模式管理、任务分发与状态联动；执行侧驱动机械臂自动调整照射方向，并通过SPI总线控制WS2812B输出跟踪、阴影、空闲、异常等状态化光效与亮度、色温变化。系统同时集成坐姿提醒功能，在提升局部照明均匀性的同时兼顾长时间学习与工作过程中的舒适性和交互性。作品充分体现了RDK X5在端侧视觉智能、嵌入式机器人协同控制和智能照明应用中的综合能力，可为桌面学习、实验工位和精细作业场景提供稳定、可扩展的智能补光方案。

\section*{第一部分 作品概述}

\subsection*{1.1 功能与特性}

本系统围绕“发现阴影、主动补光、便捷交互、健康辅助”四个目标展开设计。用户在书写、绘图或桌面操作过程中，D435会持续感知作业区域状态；当手部遮挡导致目标区域亮度下降时，系统自动判定需要补光，并驱动OpenMANIPULATOR-X调整灯具指向，使照明中心向阴影区域偏转，从而减轻局部遮挡对阅读和操作的影响。

除自动补光外，系统还提供完整的离线语音交互能力。用户可通过语音完成跟踪开启与关闭、灯光开启与关闭、亮度升高与降低、最大亮度、最小亮度、中等亮度以及暖光/冷光模式切换，避免在学习或实验过程中频繁手动触碰设备。WS2812B灯光可根据跟踪、阴影、空闲、异常等状态呈现不同光色与亮度表现，使系统工作状态更加直观。针对长时间伏案场景，系统还集成姿态提醒功能，可对肩部倾斜、身体偏移、头部过近等不良坐姿进行识别与语音提醒，从而形成“主动照明 + 语音控制 + 姿态关注”的复合功能体验。

\subsection*{1.2 应用领域}

智能无影灯系统适用于以桌面局部照明质量为核心需求的多类场景。在教学与学习环境中，该系统可服务于学生宿舍、自习室、图书馆阅览桌等场所，针对书写、绘图、阅读等任务自动减轻手部遮挡带来的阴影干扰，提升视野清晰度与长时间学习舒适度。在实验与创客环境中，系统可用于电子焊接、小型装配、模型制作、实验记录等精细作业场景，通过机械臂调整光照方向，提高目标区域的可见性与操作稳定性。

与此同时，本作品也适合应用于家庭书房、办公工位以及需要人机交互演示的智能硬件展示场景。其优势在于无需依赖云端网络即可在板端完成视觉分析、语音控制和灯光执行，既满足了嵌入式智能设备对实时性和稳定性的要求，也具备良好的展示性与扩展性，适合作为智慧照明、学习辅助和机器人协同控制方向的综合性示范作品。

\subsection*{1.3 主要技术特点}

本作品在技术实现上采用“感知层、协调层、执行层”三层板端架构。感知层以D435为视觉输入，依托RDK X5的8核ARM Cortex-A55处理器与10 TOPS BPU异构算力资源完成手部、阴影和姿态相关状态提取，其中视觉推理任务部署在BPU侧，以提升端侧感知处理效率；协调层基于ROS2 Humble组织系统消息流、模式管理与任务调度，实现视觉结果、语音指令和设备状态的统一编排；执行层由机械臂调姿与可编程灯光两部分构成，分别负责照明方向修正和光照效果输出。

在具体实现上，系统形成了清晰的板端控制闭环。视觉节点输出阴影位置、补光需求和姿态状态后，协调节点综合当前模式生成机械臂调节指令与灯光控制指令；OpenMANIPULATOR-X据此完成照明方向修正，构成“视觉感知到机械补光”的执行链路。灯光节点订阅\texttt{/light/command}，通过\texttt{spidev}在SPI \texttt{bus=1, device=0}上驱动WS2812B，并将\texttt{warm\_light\_mode}、\texttt{cool\_light\_mode}、\texttt{tracking}、\texttt{shadow}、\texttt{error}、\texttt{idle}等模式映射为对应的光色和亮度表现。离线语音模块支持亮度调节、冷暖光切换和跟踪启停，姿态提醒模块则与照明控制并行运行，体现出系统在RDK X5板端的多任务协同实现特性。

\subsection*{1.4 主要性能指标}

\begin{table}[htbp]
  \centering
  \caption{系统主要性能指标}
  \renewcommand{\arraystretch}{1.25}
  \small
  \begin{tabular}{>{\raggedright\arraybackslash}p{0.23\textwidth} >{\raggedright\arraybackslash}p{0.67\textwidth}}
    \toprule
    指标项目 & 主要能力指标 \\
    \midrule
    端侧一体化运行能力 & 系统完整运行于RDK X5平台，覆盖视觉感知、语音交互、协调控制、机械臂执行和灯光输出，满足板端嵌入式集成展示要求 \\
    视觉感知与推理能力 & 基于D435获取桌面图像，结合RDK X5 BPU完成手部、阴影和姿态相关状态提取，可在板端持续输出补光判断依据 \\
    阴影闭环调节能力 & 能够针对手部遮挡引起的局部阴影生成补光策略，并通过机械臂调姿完成照明方向修正，形成感知到执行的自动闭环 \\
    语音交互覆盖度 & 支持11类离线语音控制指令，覆盖跟踪启停、灯光启停、亮度升降、最大/最小/中等亮度以及暖光/冷光切换 \\
    灯光执行与调节能力 & 通过\texttt{spidev}在SPI \texttt{bus=1, device=0}上驱动WS2812B，支持亮度调节、冷暖光输出和状态化光效反馈 \\
    状态反馈能力 & 支持\texttt{tracking}、\texttt{shadow}、\texttt{error}、\texttt{idle}及冷暖光模式的可视化映射，增强设备运行状态的可辨识性 \\
    系统协同能力 & 基于ROS2完成\texttt{/vision/state}、\texttt{/voice/command}、\texttt{/arm/command}、\texttt{/light/command}等链路联动，支持多模块并行协同工作 \\
    健康辅助能力 & 支持肩部倾斜、身体偏移、头部过近等姿态异常提醒，使照明系统兼具学习过程中的人机辅助功能 \\
    \bottomrule
  \end{tabular}
\end{table}

\subsection*{1.5 主要创新点}

\begin{enumerate}[label=\arabic*.]
  \item 构建了面向桌面学习场景的全板端智能无影灯方案，将视觉推理、语音交互、机械臂控制和灯光驱动统一部署在RDK X5平台上，突出嵌入式一体化实现能力。
  \item 将RDK X5 BPU视觉推理与机械臂定向补光闭环结合，使系统能够从“看见阴影”进一步发展为“理解阴影并主动修正光照方向”，提升了智能照明系统的主动性。
  \item 将WS2812B可编程灯光经SPI纳入系统状态反馈链路，不仅实现冷暖光与亮度调节，还通过\texttt{tracking}、\texttt{shadow}、\texttt{error}、\texttt{idle}等模式提升设备运行状态的可感知性。
  \item 在照明辅助之外融入姿态提醒功能，使作品从单一照明设备拓展为兼顾学习健康管理的人机协同系统，增强了竞赛作品的场景完整性与应用价值。
\end{enumerate}

\subsection*{1.6 设计流程}

本作品的工程设计流程遵循“需求定义、方案分解、板端部署、联调验证、场景评估”的推进路径。首先围绕桌面书写与精细操作中常见的遮挡阴影问题，明确自动补光、语音控制、姿态提醒和状态反馈等核心需求，并据此确定系统评价重点；随后完成视觉、语音、协调控制、机械臂和灯光等模块的接口划分与数据流设计，建立统一的ROS2通信链路；在实现阶段，将视觉推理部署到RDK X5 BPU，在板端完成离线语音控制与WS2812B的SPI驱动集成，并对机械臂调姿策略和灯光模式映射进行参数联调；最后通过阴影跟踪、亮度与冷暖光切换、姿态提醒和状态灯联动等典型场景进行整机验证，确保系统具备稳定、连贯的板端运行效果。

% \begin{figure}[htbp]
%   \centering
%   \includegraphics[width=0.82\textwidth]{figures/design_flow.png}
%   \caption{智能无影灯系统设计流程示意图}
% \end{figure}
\ImagePlaceholder{智能无影灯系统设计流程示意图}{4.5cm}

\section*{第二部分 系统组成及功能说明}

\subsection*{2.1 整体介绍}

本系统围绕RDK X5构建完整的板端智能照明闭环，整体由视觉感知、语音交互、协调决策、机械臂执行、灯光输出和姿态提醒六个功能单元组成。D435持续采集桌面作业区域图像与深度信息，板载麦克风负责接收唤醒词与离线控制指令，RDK X5在ROS2 Humble框架下完成多源信息汇聚、任务调度与设备控制，OpenMANIPULATOR-X负责调整灯体指向，WS2812B灯光单元负责输出照明效果与系统状态反馈。整机形成了“感知、决策、执行、反馈”一致化运行路径，能够面向书写、阅读和精细操作场景实现自动补光与自然交互。

从数据流角度看，视觉模块向\texttt{/vision/state}发布手部位置、阴影状态、补光需求、姿态评价以及深度定位结果，语音模块向\texttt{/voice/command}发布识别后的控制指令；\texttt{system\_coordinator}根据当前模式和环境状态生成系统模式消息\texttt{/system/mode}，并进一步生成机械臂控制量与灯光控制量，分别发布到\texttt{/arm/command}和\texttt{/light/command}；\texttt{arm\_control}再把方向性调节量映射为MoveIt Servo可执行的\texttt{JointJog}速度命令，最终发送到\texttt{/servo\_node/delta\_joint\_cmds}。与此同时，姿态提醒链路并行运行，在检测到肩部倾斜、身体左右偏移或头部距离桌面过近时给出语音反馈，使系统不仅具备无影补光能力，也具备学习辅助能力。

\subsection*{2.2 硬件系统介绍}

系统硬件以RDK X5为主控核心，围绕视觉输入、音频输入、机械执行与灯光输出四类接口展开布置。各硬件单元既服务于板端算法运行，也直接决定整机的交互质量与动作响应能力，因此在结构设计上采用“主控集中、传感分布、执行分层”的组织方式。

\subsubsection*{2.2.1 硬件整体介绍}

RDK X5承担整机的嵌入式计算、ROS2节点运行、BPU视觉推理调度和外设控制，是系统的统一控制中心。视觉感知部分采用Intel RealSense D435，利用其彩色图像与深度信息同时支持手部检测、阴影区域定位、补光目标估计和三维手点定位；音频交互部分采用麦克风采集用户语音，用于触发离线唤醒与命令识别；执行部分采用OpenMANIPULATOR-X机械臂搭载灯体，通过关节联动改变照明方向；灯光输出部分采用WS2812B可编程RGB光源，既承担状态化光效显示，也承担亮度和色温模式切换的结果呈现。

从接口连接关系看，D435与麦克风均作为RDK X5的输入设备接入，机械臂与灯光单元分别作为执行设备接收主控下发的控制命令。前者主要负责空间指向调整，后者负责光色、亮度和状态反馈，两者共同构成补光系统的物理输出端。因此，RDK X5并非仅作为算法开发平台，而是作为整套智能无影灯的板端核心，直接完成从感知到动作和灯光反馈的全过程控制。

% \begin{figure}[htbp]
%   \centering
%   \includegraphics[width=0.80\textwidth]{figures/system_whole_machine.jpg}
%   \caption{智能无影灯整机实物图}
% \end{figure}
\ImagePlaceholder{智能无影灯整机实物图}{5.0cm}

% \begin{figure}[htbp]
%   \centering
%   \includegraphics[width=0.82\textwidth]{figures/rdk_x5_peripheral_connection.png}
%   \caption{RDK X5与外设连接关系图}
% \end{figure}
\ImagePlaceholder{RDK X5与外设连接关系图}{5.0cm}

\subsubsection*{2.2.2 机械设计介绍}

机械部分以OpenMANIPULATOR-X为主体，并将末端执行器由传统夹持用途转化为灯体安装与定向补光载体。该设计使灯光不再是固定角度照明，而是能够随着桌面阴影分布动态调整照射方向。机械结构重点利用底座回转方向与前级俯仰方向形成对照明区域的主调节自由度，使末端灯体能够围绕工作台面完成水平偏转与俯仰修正，从而覆盖书写、阅读、实验等典型近距离场景。

在当前实现中，系统控制上主要使用\texttt{joint1}和\texttt{joint2}作为照明方向调节代理量，对应水平偏航与垂向俯仰两类关键动作。这样的机械设计具有两个工程优势：一是控制链路清晰，视觉偏差可以直接映射为灯体方向修正量；二是结构稳定，机械臂本体、灯体和观察区域之间能够维持稳定的几何关系，便于后续继续扩展更精细的末端姿态控制或灯具结构优化。

% \begin{figure}[htbp]
%   \centering
%   \includegraphics[width=0.76\textwidth]{figures/arm_lamp_structure.jpg}
%   \caption{机械臂与灯体结构图}
% \end{figure}
\ImagePlaceholder{机械臂与灯体结构图}{4.8cm}

\subsubsection*{2.2.3 电路各模块介绍}

电路系统围绕“主控接口层、传感接入层、执行驱动层、灯光输出层”四部分组织。主控接口层以RDK X5的USB、串口和SPI资源为中心，对外连接视觉、语音、机械臂和灯光模块。D435通过高速数据接口向主控传输彩色图像与深度数据，麦克风模块向语音识别链路提供音频输入，机械臂控制链路通过USB串口形式与主控通信，用于下发运动控制命令并接入底层驱动。

灯光输出层采用WS2812B作为可编程光源，控制路径由RDK X5的\texttt{SPI}接口直接驱动，主控向灯光单元发送编码后的时序数据以完成亮度、色温和状态光效输出。这一路径已经在板端实现，不依赖额外上位机控制；从接口拓扑上看，WS2812B与RDK X5之间形成了明确的主控到光源直连链路，与视觉输入链路、语音输入链路和机械臂串口控制链路共同构成整机电路连接关系。整体上，电路设计强调接口明确、链路闭合和板端可落地性，符合嵌入式智能照明系统的工程实现要求。

% \begin{figure}[htbp]
%   \centering
%   \includegraphics[width=0.76\textwidth]{figures/ws2812_spi_connection.jpg}
%   \caption{WS2812B SPI驱动连接与灯效展示}
% \end{figure}
\ImagePlaceholder{WS2812B SPI驱动连接与灯效展示}{4.8cm}

\subsection*{2.3 软件系统介绍}

软件系统采用基于ROS2 Humble的模块化分层架构，在RDK X5板端完成感知、决策、控制与反馈的统一部署。整个软件栈既包含消息接口定义与节点通信机制，也包含BPU视觉推理、离线语音解析、机械臂速度控制和灯光驱动等具体实现模块，是支撑整机稳定运行的核心部分。

\subsubsection*{2.3.1 软件整体介绍}

软件平台运行于RDK X5的Ubuntu 22.04环境，采用ROS2 Humble作为通信中间件，并以\texttt{shadow\_lamp\_interfaces}定义统一消息接口。系统通过话题发布与订阅的方式组织多节点协作，其中视觉模块负责连续输出环境状态，语音模块负责输出用户意图，协调模块负责把感知结果转化为控制命令，机械臂与灯光模块负责执行动作与反馈结果。由于各模块围绕统一接口解耦，系统能够在板端实现并行运行与独立调试，具备较好的工程扩展性。

在算力分配上，视觉侧重点利用RDK X5的\texttt{BPU}资源完成手部、阴影和姿态相关状态提取，CPU侧则承担ROS2节点调度、消息编排、语音解析和设备控制等任务。这种“BPU负责高效感知、CPU负责系统控制”的协同方式，使系统既能够保持端侧智能分析能力，又能保证命令链路简洁稳定。结合\texttt{/vision/state}、\texttt{/voice/command}、\texttt{/system/mode}、\texttt{/arm/command}、\texttt{/light/command}和\texttt{/servo\_node/delta\_joint\_cmds}等关键链路，软件系统形成了清晰的板端实时控制闭环。

% \begin{figure}[htbp]
%   \centering
%   \includegraphics[width=0.82\textwidth]{figures/software_architecture.png}
%   \caption{软件架构与运行界面示意图}
% \end{figure}
\ImagePlaceholder{软件架构与运行界面示意图}{5.0cm}

\subsubsection*{2.3.2 软件各模块介绍}

\begin{enumerate}[label=\arabic*.]
  \item \textbf{视觉感知模块。}视觉模块负责生成系统对桌面工作区的基础理解，并向\texttt{/vision/state}发布统一状态消息。其核心输出包括\texttt{hand\_detected}、\texttt{hand\_center\_x/y}、\texttt{shadow\_detected}、\texttt{shadow\_center\_x/y}、\texttt{shadow\_area}、\texttt{needs\_relight}、\texttt{suggested\_target\_x/y}，以及姿态与深度相关字段\texttt{posture\_ok}、\texttt{posture\_issue}、\texttt{posture\_score}、\texttt{hand\_depth\_valid}、\texttt{hand\_depth\_m}、\texttt{hand\_point\_camera\_x/y/z}。这些信息一方面用于二维阴影补光，另一方面为三维指向控制和姿态提醒提供输入依据。

% \begin{figure}[htbp]
%   \centering
%   \includegraphics[width=0.74\textwidth]{figures/hand_shadow_detection_result.jpg}
%   \caption{手部与阴影检测结果示意图}
% \end{figure}
\ImagePlaceholder{手部与阴影检测结果示意图}{4.6cm}

  \item \textbf{语音交互模块。}语音模块围绕离线唤醒和固定命令识别展开工作，通过唤醒词触发控制流程，并将识别结果发布到\texttt{/voice/command}。当前已实现的命令覆盖跟踪开启与关闭、灯光开启与关闭、亮度升高与降低、最大亮度、最小亮度、中等亮度以及暖光/冷光模式切换，可满足桌面使用过程中的高频操作需求。该模块的意义在于将“用户意图”转化为系统模式切换条件，使整套设备在不依赖网络的情况下仍具备自然交互能力。

% \begin{figure}[htbp]
%   \centering
%   \includegraphics[width=0.76\textwidth]{figures/voice_control_flow_interface.png}
%   \caption{离线语音控制流程与界面示意图}
% \end{figure}
\ImagePlaceholder{离线语音控制流程与界面示意图}{4.6cm}

  \item \textbf{协调决策模块。}协调模块是整机的软件中枢，负责融合视觉状态、语音命令和当前运行模式，并生成\texttt{SystemMode}、\texttt{ArmCommand}与\texttt{LightCommand}。在阴影跟踪场景下，其核心计算遵循手影偏差归一化思想，可表示为
  \[
    dx = x_{shadow} - x_{hand}, \qquad dy = y_{shadow} - y_{hand}
  \]
  \[
    yaw = \frac{dx}{W/4}, \qquad pitch = \frac{dy}{H/4}
  \]
  其中，$W$和$H$分别表示输入图像的宽度与高度。再结合阈值抑制和幅值限制生成机械臂方向控制量。模块支持\texttt{shadow\_follow}、\texttt{hand\_follow}、\texttt{hand\_point\_3d}、\texttt{hand\_point\_3d\_yaw\_only}等多种工作风格，其中三维模式会综合\texttt{hand\_depth\_valid}、\texttt{hand\_depth\_m}和\texttt{hand\_point\_camera\_x/y/z}等信息输出空间目标。最终形成的\texttt{SystemMode}消息通过\texttt{/system/mode}发布，用于表征当前模式与跟踪开关状态；同时生成的\texttt{ArmCommand}包含\texttt{mode}、\texttt{follow\_enabled}、\texttt{target\_x/y/z}、\texttt{yaw}、\texttt{pitch}等字段，并通过\texttt{/arm/command}发布；灯光相关控制则通过\texttt{/light/command}下发。

  \item \textbf{机械臂控制模块。}机械臂控制模块位于上层决策与底层Servo之间，负责把\texttt{/arm/command}中的偏航、俯仰或目标跟随意图转化为实际关节速度命令。当前实现中，\texttt{arm\_control}主要将\texttt{yaw}映射为\texttt{joint1}速度，将\texttt{pitch}映射为\texttt{joint2}速度，并封装为MoveIt Servo使用的\texttt{JointJog}消息后发布到\texttt{/servo\_node/delta\_joint\_cmds}。这一链路保证了协调层生成的是面向任务的抽象控制量，而执行层接收的是可直接作用于OpenMANIPULATOR-X的底层速度命令。

  \item \textbf{灯光控制模块。}灯光控制模块订阅\texttt{/light/command}，并根据\texttt{mode}、\texttt{brightness}、\texttt{color\_temperature}和\texttt{enabled}等字段驱动WS2812B输出。板端实现中，\texttt{light\_control/ws2812\_spi.py}通过\texttt{spidev}完成WS2812B帧编码与SPI发送，属于已实现的正式控制路径。\texttt{light\_controller\_node}将\texttt{warm\_light\_mode}、\texttt{cool\_light\_mode}、\texttt{tracking}、\texttt{shadow}、\texttt{error}、\texttt{idle}等状态映射为不同RGB表现，使灯光既是照明执行部件，也是系统状态输出部件。

  \item \textbf{姿态提醒模块。}姿态提醒与补光控制并行工作，主要利用视觉输出中的\texttt{posture\_ok}、\texttt{posture\_issue}和\texttt{posture\_score}判断用户坐姿是否稳定。当出现肩部倾斜、身体左右偏移或头部距离桌面过近等问题并持续满足触发条件时，系统会调用本地语音反馈链路给出提醒。该模块使智能无影灯从单一照明设备扩展为兼顾健康辅助的人机协同系统，提高了作品在长期学习场景中的实用性。
\end{enumerate}

\section*{第三部分 完成情况及性能参数}

\subsection*{3.1 整体介绍}

基于RDK X5的智能无影灯系统目前已完成整机联调，形成可直接演示的板端样机。整机由RDK X5主控、D435视觉模组、离线语音输入链路、OpenMANIPULATOR-X机械臂灯体以及基于SPI驱动的WS2812B光源组成，在ROS2 Humble框架下实现视觉感知、语音交互、机械调姿、灯光反馈与姿态提醒的一体化运行。视觉状态、语音指令、机械臂控制量和灯光控制量均在板端生成并执行，构成完整的嵌入式闭环。

当前系统已经具备明确、可复现的运行参数：语音控制覆盖11类主命令，姿态提醒触发持续时间为4~s、同类问题提醒冷却时间为45~s；阴影跟踪采用15~px死区抑制抖动，在$640\times480$参考图像坐标下生成归一化的偏航与俯仰调节量；机械臂控制链路默认速度系数为0.6；WS2812B灯控采用\texttt{SPI bus=1, device=0, max\_speed\_hz=2400000}配置。整机能够围绕桌面书写与精细作业场景完成自动补光、亮度调节、冷暖光切换和姿态提醒等核心功能。

\subsection*{3.2 工程成果}

围绕“机械可安装、电路可连接、软件可运行、灯光可执行”的工程目标，项目已完成整机样机搭建、板端链路打通和关键参数固化。当前成果覆盖机械结构、接口连接、BPU视觉推理、离线语音交互、机械臂执行和WS2812B灯控执行，能够支撑完整的竞赛展示与后续优化。

\subsubsection*{3.2.1 机械成果}

机械方面，项目完成了以OpenMANIPULATOR-X为主体的智能无影灯样机搭建，并将灯体稳定安装到机械臂末端，使机械臂由传统抓取执行器转变为定向补光执行器。整机可围绕桌面目标区域完成水平偏转与俯仰修正，满足书写、阅读和小型实验操作场景中对局部照射角度调整的需求。当前控制链路以\texttt{joint1}和\texttt{joint2}作为主要调节自由度，分别承担水平方向偏航和垂直方向俯仰控制。

机械执行结果已经形成稳定参数闭环：阴影跟踪在$640\times480$参考图像坐标下计算偏差，15~px以内采用死区保持，超出死区后将偏差归一化为$[-1,1]$范围内的\texttt{yaw}/\texttt{pitch}控制量；\texttt{arm\_control}再以0.6的速度系数映射为\texttt{JointJog}关节速度命令并发送到\texttt{/servo\_node/delta\_joint\_cmds}。这一结果说明系统已经实现从视觉偏差到灯体动作的可见机械响应。

\subsubsection*{3.2.2 电路成果}

电路方面，项目已经完成以RDK X5为中心的主控与外设连接方案，形成了视觉输入、语音输入、机械执行和灯光输出的完整接口体系。D435持续向主控输入彩色图像与深度信息，麦克风负责采集离线语音控制音频，机械臂通过USB串口链路接收控制命令，WS2812B灯光模块由RDK X5的SPI接口直接驱动，整套设备在板端完成接入、消息传递和执行联动。

其中，WS2812B的SPI驱动链路已经完成工程落地，是电路成果中的重点。软件侧通过\texttt{spidev}在\texttt{bus=1, device=0}、\texttt{max\_speed\_hz=2400000}配置下发送编码后的灯控数据流，实现对WS2812B的亮度、色温与状态光效控制。该结果表明RDK X5不仅承担算法与控制逻辑，也承担板端外设驱动和光源执行任务，整机接口关系明确、执行链路闭合。

\subsubsection*{3.2.3 软件成果}

软件方面，项目已经完成围绕ROS2 Humble构建的板端多节点系统，核心软件成果包括统一消息接口、BPU视觉感知链路、离线语音识别链路、协调决策链路、机械臂Servo控制链路以及灯光控制链路。\texttt{shadow\_lamp\_interfaces}完成了视觉、语音、模式、机械臂和灯光消息定义，\texttt{vision\_perception}负责输出手部、阴影、姿态和深度相关状态，\texttt{system\_coordinator}负责生成控制决策，\texttt{arm\_control}负责把方向调节量映射为\texttt{JointJog}消息，\texttt{light\_control}负责落实灯光执行。

当前软件链路已经形成明确参数化结果。视觉侧利用RDK X5的BPU完成手部、阴影和姿态相关状态提取，CPU侧承担ROS2消息调度、模式管理、语音解析和设备控制；\texttt{/vision/state}消息提供19个业务字段和1个时间戳字段，覆盖二维手影信息、姿态评分以及三维手点坐标；语音模块支持11类主命令，覆盖跟踪启停、灯光启停、亮度升降、最大/最小/中等亮度和暖光/冷光切换；协调模块已实现\texttt{shadow\_follow}、\texttt{hand\_follow}、\texttt{hand\_point\_3d}和\texttt{hand\_point\_3d\_yaw\_only}等控制模式。软件系统已经能够稳定支撑整机联动运行。

\subsubsection*{3.2.4 灯控成果}

灯控方面，项目已经完成从上层控制命令到WS2812B实际输出的完整实现，形成了既能照明调节、又能状态提示的灯光执行链路。系统通过\texttt{/light/command}接收\texttt{mode}、\texttt{brightness}、\texttt{color\_temperature}、\texttt{enabled}等控制字段，在板端映射为WS2812B对应的RGB数据并通过SPI发送。用户可通过离线语音完成灯光开启与关闭、亮度升高与降低、最大亮度、最小亮度、中等亮度以及暖光/冷光模式切换，说明亮度与色温控制已形成实际运行结果。

当前灯控系统已实现6类命名模式输出，分别为\texttt{warm\_light\_mode}、\texttt{cool\_light\_mode}、\texttt{tracking}、\texttt{shadow}、\texttt{error}和\texttt{idle}。其中暖光与冷光用于照明模式切换，\texttt{tracking}、\texttt{shadow}、\texttt{error}和\texttt{idle}用于运行状态指示。该设计使评审者和用户无需查看日志即可直接感知系统所处状态，提升了整机的可视化反馈效果和现场展示效果。

\subsection*{3.3 特性成果}

综合整机联调结果，本作品已经形成较完整的竞赛型技术特征：一是所有关键链路均围绕RDK X5板端组织，具备端侧一体化实现能力；二是BPU视觉感知、机械臂调姿和WS2812B灯控之间形成闭环，使系统具备自动减弱桌面阴影的能力；三是离线语音交互直接作用于跟踪与灯控，提升了使用便捷性；四是姿态提醒与状态灯反馈增强了作品的场景完整性与可展示性。相关运行参数如表所示。

\begin{table}[htbp]
  \centering
  \caption{整机完成情况与关键性能参数}
  \renewcommand{\arraystretch}{1.25}
  \small
  \begin{tabular}{>{\raggedright\arraybackslash}p{0.25\textwidth} >{\raggedright\arraybackslash}p{0.65\textwidth}}
    \toprule
    项目 & 完成结果与参数说明 \\
    \midrule
    板端主控平台 & RDK X5，Ubuntu 22.04，ROS2 Humble；视觉感知任务部署在BPU侧，系统控制与外设驱动运行在CPU侧 \\
    视觉控制参数 & 参考图像坐标为$640\times480$，阴影跟踪死区为15~px，\texttt{yaw}/\texttt{pitch}输出归一化并限制在$[-1,1]$范围 \\
    视觉状态输出 & \texttt{/vision/state}提供19个业务字段和1个时间戳字段，覆盖手部、阴影、姿态、深度和三维手点坐标信息 \\
    阴影补光闭环 & 已实现\texttt{/vision/state} $\rightarrow$ \texttt{/arm/command} $\rightarrow$ \texttt{/servo\_node/delta\_joint\_cmds}控制链路 \\
    机械调节参数 & 基于\texttt{joint1}与\texttt{joint2}的偏航/俯仰调节，默认速度系数\texttt{velocity\_scale=0.6} \\
    灯光控制参数 & 已完成WS2812B的SPI驱动，运行参数为\texttt{bus=1, device=0, max\_speed\_hz=2400000} \\
    语音控制能力 & 已完成11类离线控制指令，覆盖跟踪启停、灯光启停、亮度升降、最大/最小/中等亮度及冷暖光切换 \\
    状态反馈能力 & 已实现6类命名模式：\texttt{warm\_light\_mode}、\texttt{cool\_light\_mode}、\texttt{tracking}、\texttt{shadow}、\texttt{error}、\texttt{idle} \\
    健康辅助能力 & 已支持肩部倾斜、身体偏移、头部过近3类姿态提醒，触发持续时间4~s、同类问题提醒冷却时间45~s \\
    \bottomrule
  \end{tabular}
\end{table}

% \begin{figure}[htbp]
%   \centering
%   \includegraphics[width=0.82\textwidth]{figures/shadow_reduction_effect.jpg}
%   \caption{阴影减弱与控制效果展示}
% \end{figure}
\ImagePlaceholder{阴影减弱与控制效果展示}{5.0cm}

\section*{第四部分 总结}

\subsection*{4.1 可扩展之处}

本作品已经完成板端一体化原型，但仍有明确且现实的扩展方向。首先，在视觉侧可以继续强化RDK X5 BPU上的模型能力，例如引入更稳定的手部关键点模型、桌面区域分割模型或阴影检测模型，提高复杂背景和不同光照条件下的识别鲁棒性；同时结合深度信息与外参标定，进一步把当前以\texttt{joint1}/\texttt{joint2}为主的方向调节升级为更精细的末端照明位姿控制，使灯具姿态、照射中心和照度分布能够按三维目标进行优化。

其次，在灯控和交互侧还可以继续完善面向实际使用的功能。灯光部分可扩展为阅读模式、专注模式、夜间模式、演示模式等更细分的照明场景，并进一步引入亮度渐变、区域分组或照度闭环调节；语音部分可继续增加本地离线命令词，如系统状态查询、模式切换、定时休息提醒和一键演示等；协调层则可继续增加安全限位、异常自检、空闲待机和使用统计功能。上述扩展都可以沿用当前的RDK X5板端架构和ROS2接口体系推进，工程连续性较强。

\subsection*{4.2 心得体会}

本项目的实施过程表明，智能无影灯并不是单一算法或单一硬件即可完成的作品，而是一个典型的多模块协同系统。真正决定整机完成度的，除了视觉识别准确性，还包括消息接口是否统一、机械动作是否直观、灯光反馈是否清晰、语音交互是否自然，以及这些模块在板端是否能够稳定协同。RDK X5在本项目中承担了视觉推理、外设驱动、控制编排和整机联动的核心职责，其板端系统集成能力是作品得以落地的关键。

在工程实现层面，我们也更清楚地认识到“感知、决策、执行、反馈”闭环设计的重要性。将BPU感知结果通过ROS2消息规范化之后，机械臂控制和灯控逻辑的联调效率明显提升；把WS2812B灯光纳入系统反馈链路之后，整机状态的可观察性也显著增强。同时也需要看到当前方案的工程边界：现阶段机械控制主要依赖\texttt{joint1}/\texttt{joint2}作为灯体方向代理量，已能够满足阴影跟踪演示需求，但距离更精细的末端灯头位姿控制、照射中心优化和照度闭环调节仍有提升空间。总体来看，本项目积累了端侧智能照明、机器人协同控制和嵌入式系统整合方面的直接经验，也为后续继续打磨成更成熟的产品化原型提供了清晰方向。

\section*{第五部分 参考文献}

\begin{thebibliography}{99}
  \bibitem{rdkx5}
  地平线机器人. RDK X5开发者文档[EB/OL]. \url{https://developer.d-robotics.cc/rdkx5}, 2026.

  \bibitem{rdkmodelzoo}
  D-Robotics. RDK Model Zoo and BPU Deployment Examples[EB/OL]. \url{https://github.com/D-Robotics/rdk_model_zoo}, 2026.

  \bibitem{ros2}
  Open Robotics. ROS 2 Documentation: Humble Hawksbill[EB/OL]. \url{https://docs.ros.org/en/humble/}, 2026.

  \bibitem{realsense}
  Intel Corporation. Intel RealSense Depth Camera D435 Datasheet and Developer Documentation[EB/OL]. \url{https://dev.intelrealsense.com/docs}, 2026.

  \bibitem{openmanipulator}
  ROBOTIS. OpenMANIPULATOR-X e-Manual[EB/OL]. \url{https://emanual.robotis.com/docs/en/platform/openmanipulator_x/overview/}, 2026.

  \bibitem{moveitservo}
  PickNik Robotics. MoveIt Documentation: Realtime Arm Servoing[EB/OL]. \url{https://moveit.picknik.ai/humble/doc/examples/realtime_servo/realtime_servo_tutorial.html}, 2026.

  \bibitem{ws2812b}
  Worldsemi. WS2812B LED Datasheet[EB/OL]. \url{https://cdn.sparkfun.com/assets/e/6/1/f/4/WS2812B-LED-datasheet.pdf}, 2026.

  \bibitem{vosk}
  Vosk Team. Vosk Offline Speech Recognition Toolkit[EB/OL]. \url{https://alphacephei.com/vosk/}, 2026.

  \bibitem{mediapipe}
  Google. MediaPipe Solutions Guide[EB/OL]. \url{https://developers.google.com/mediapipe}, 2026.
\end{thebibliography}

\end{document}
